\documentclass[10pt,a4paper]{amsart}
\usepackage[margin=19mm]{geometry}
\usepackage{amsmath,amssymb,mathtools}
\usepackage[scr=rsfs,cal=euler]{mathalfa}
\usepackage{xfrac}
\usepackage[unicode,hidelinks,bookmarks=false]{hyperref}
\newcommand{\ie}{i.e.\ }
\newcommand{\cf}{cf.\ }
\newcommand{\resp}{respectively\ }
\newcommand{\Subsection}{\subsection}
\providecommand{\nobreakdash}{\nobreak\hbox{-}}
\setlength{\emergencystretch}{2em}
\multlinegap0pt
\usepackage{amssymb}
\newtheorem{prop}{Proposition}
\newtheorem{lem}{Lemma}
\newcommand{\A}{{\mathcal A}}
\newcommand{\al}{{\alpha}}
\newcommand{\be}{{\beta}}
\title{On conic-line arrangements with nodes, tacnodes, and ordinary triple points}
\author{Alexandru Dimca, Piotr Pokora$^{\ast}$}
\date{Source: arXiv:2111.12349v3 (CC BY 4.0)}
\begin{document}
\maketitle

\noindent\emph{Source and licence.} On conic-line arrangements with nodes, tacnodes, and ordinary triple points. Version arXiv:2111.12349v3 (CC BY 4.0), distributed by arXiv under the Creative Commons Attribution 4.0 International licence, http://creativecommons.org/licenses/by/4.0/, as recorded in the arXiv licence metadata for this exact version and shown on the abstract page https://arxiv.org/abs/2111.12349v3. Full licence notice: \texttt{licenses/arxiv-2111.12349-CC-BY-4.0.txt}.\par\smallskip

\noindent\textit{Editorial scope.} Counted unit: the article's prop prop11b (source0.tex lines 653--656). The article's complete original proof of it is delivered with it (source0.tex lines 657--733, one \texttt{proof} environment of 77 lines). No mathematics is written, repaired or re-derived: the statement and the proof are the article's own lines, in the article's own notation, and no retained line is substituted or re-flowed.  Retained uncounted as the source-local context the proof needs: lines 333--338; lines 455--458; lines 460--463; lines 489--493; lines 577--580; lines 582--585.  Nothing was omitted from any retained span, so the package carries no omission record.  No retained line contains a citation, so the wrapper carries no bibliography and no bibliography entry.  No retained line contains a figure environment, an image-inclusion command or drawing code, so no image is delivered and the package declares no asset dependency.  Editorial formatting only, in addition: an AMS \texttt{amsart} preamble that restates the article's own declarations and macros used by the retained lines, namely the environments \texttt{prop}, \texttt{lem} on separate counters, together with the standard packages those lines need.  The retained environments print in this document as Proposition 1, Lemma 1 in order of appearance, whereas the article prints the same items with the numbering of its own full document.  The complete native source of the article as distributed in the arXiv e-print for this exact version is bundled unchanged as \texttt{sources/118826/source0.tex}.
\begin{prop}
\label{prop:bound9}
Let $C \, : \, f=0$ be a reduced curve of degree $m$ in $\mathbb{P}^{2}_{\mathbb{C}}$ having only nodes, tacnodes, and ordinary triple points as singularities. Then
$${\rm mdr}(f) \geq \frac{2}{3}m-2.$$
In particular, if $C$ is free, then $m \leq 9$.
\end{prop}

\begin{equation}
\label{E7"}
n_2+n_3 \leq \frac{3}{2}(d-1).
\end{equation}

\begin{equation}
\label{E7'}
t \geq k^2 - k + kd +\frac{(d-1)(d-3)}{4}-n_3.
\end{equation}

\begin{lem}
\label{lem16}
Let $C_1$ and $C_2$ be two smooth conics in $\mathcal{CL}$ such that 
$|C_1 \cap C_2|=2$. Then any line $L$ in $\mathcal{CL}$ is tangent to at most one of the conics $C_1$ and $C_2$.
\end{lem}

\begin{equation}
\label{E11a}
n_2+ n_3 \leq 4
\end{equation}

\begin{equation}
\label{E11b}
t \geq k^2-k+4k + \frac{3}{4} - n_3.
\end{equation}

\begin{prop}
\label{prop11b}
Let $\mathcal{CL}$ be an arrangement of $d=5$ lines and $k \geq 1$ smooth conics having only nodes, tacnodes, and ordinary triple points. Then $\mathcal{CL}$  is never free.
\end{prop}

\begin{proof}
If we assume that $\mathcal{CL}$ is free. Using Proposition \ref{prop:bound9}, we get $m = 2k+5 \leq 9$, and hence $k\leq 2$.
Then formula (\ref{E7"}) implies that
\begin{equation}
\label{E11a2}
n_2+ n_3 \leq 6
\end{equation}
and by (\ref{E7'}) we get
\begin{equation}
\label{E11b2}
t \geq k^2-k+5k+2-n_3.
\end{equation}

\medskip

{\bf Case $k=2$.}

\medskip

If $m_2=1$, then any line is tangent to at most one conic, which follows from Lemma \ref{lem16}, so we have $t \leq 7$ and \eqref{E11b2} yields a contradiction.

If $m_3=1$, the inequality \eqref{E11b2} implies $t + n_3 \geq 14$.
We can use $2$ lines, say $L_1$ and $L_2$, to create $4$ new tacnodes ($5$ in total) and a new node ($3$ in total). The third line $L_3$ can produce $2$ triple points if it passes through the $2$ nodes in the intersection $C_1 \cap C_2$, but will create $2$ new nodes as intersections with $L_1$ and $L_2$. It means that the remaining $2$ lines should not create any new node or triple point, which is impossible.  If $L_3$ passes through the node $L_1\cap L_2$, then it will be secant to both $C_1$ and $C_2$, thus creating 
$4$ new nodes, a contradiction with \eqref{E11a}.
Finally, if $L_3$ is tangent to one of the conics, it would be secant for the other, creating $4$ new nodes, $2$ on that second conic and $2$ on $L_1 \cap L_2$. This contradicts \eqref{E11a}.

If $m_4=1$, $C_1\cap C_2$ gives rise to 4 nodes. As soon as we add $3$ lines, at least $3$ new nodes will be created (and some old nodes transformed in triple points), but in any case we get a contradiction with \eqref{E11a}.

Hence the case $k=2$ is impossible.
\medskip

{\bf Case: $k=1$.}

\medskip
The above shows that the only possibility to have \eqref{E11a2} and \eqref{E11b2} satisfied is that $k=1$ and $n_3 \geq 2$.
If we look on the (sub)arrangement $\mathcal{A}$ of the $5$ lines in $\mathcal{CL}$, there are $3$ cases to discuss.

\medskip

{\bf Case 1: $\mathcal{A}$ has $2$ triple points.}

\medskip

If we denote by $p_1$ and $p_2$ the two triple points, then the line determined by these points must be in $\mathcal{A}$. We denote this line by $L$. The other $2$ lines passing through $p_1$ ($p_2$, respectively) are denoted by $L_1$ and $L_1'$ ($L_2$ and $L_2'$, respectively). There are $4$ nodes in the arrangement $\mathcal{A}$, the intersections $L_1 \cap L_2$, $L_1 \cap L_2'$,
$L_1' \cap L_2$, and $L_1' \cap L_2'$.

If $n_3=2$, we get from \eqref{E11b2} that $t \geq 5$, which is impossible since $L$, $L_1$ and $L_1'$ cannot be all tangent to the conic.

If $n_3=3$, the new triple point $q_1$ is at the intersection of two secant lines with the conic. In this case, we get from \eqref{E11b} that $t \geq 4$, which is impossible, since the $2$ lines meeting at $q_1$ cannot be tangent to the conic.

If $n_3=4$, then there are $2$ triple points $q_1$ and $q_2$ situated on the conic, coming from the intersection of at least $3$ lines from $\mathcal{A}$ with the conic. In this case, we get from \eqref{E11b} that $t \geq 3$, which is impossible since the lines passing through $q_1$ or $q_2$ cannot be tangent to the conic.

If $n_3=5$, then there are $3$ triple points $q_1$, $q_2$ and $q_3$ situated on the conic, coming from the intersection of at least $4$ lines from $\mathcal{A}$ with the conic. In this case, we get from \eqref{E11b} that $t \geq 2$, which is impossible since the lines passing through $q_1$, $q_2$ or $q_3$ cannot be tangent to the conic.

Finally, if $n_3=6$, then $n_2=0$, the conic, call it $Q$, passes through all the $4$ nodes of $\mathcal{A}$ and it is tangent to the line $L$. The conics passing through the $4$ nodes form a pencil of conics, determined by the degenerate conics $Q_1$ and $Q_2$, where $Q_1$ (resp. $Q_2$) is  the union $L_1 \cup L_1'$ (resp. $L_2 \cup L_2'$).
Note that these two degenerate conics $Q_1$ and $Q_2$ meet the line $L$ in one point, namely $Q_1\cap L=p_1$ and $Q_2\cap L=p_2$.
The conic $Q$ is in this pencil $\al Q_1 + \be Q_2$, and meets
the line $L$ in one point as well. This is a contradiction, since in a pencil only two members can meet a given line in a single point. To check this claim, we assume that $L$ is given by $x=0$. Then the condition that $\al Q_1 + \be Q_2$ meets
line $L$ in one point is a quadratic form in $y,z$, with coefficients being linear forms in $\al, \be$, such that it has zero discriminant, which yields the vanishing of the quadratic form in $\al, \be$.

\medskip

{\bf Case 2: $\mathcal{A}$  has a triple point.}

\medskip

It follows that $\mathcal{A}$ has $7$ nodes. To create $n_3$ triple points in $\mathcal{CL}$, the conic passes through $n_3-1$ nodes of the line arrangement $\A$. Hence $\mathcal{CL}$ has $n_3$ triple points and at least $7-(n_3-1)$ double points. It follows that $n_2+n_3 \geq 8$, a contradiction with respect to \eqref{E11a2}.

\medskip

{\bf Case 3: $\mathcal{A}$ is a nodal arrangement.}

\medskip

It follows that now  $\mathcal{A}$ has $10$ nodes. To create $n_3$ triple points in $\mathcal{CL}$, the conic passes through $n_3$ nodes of the line arrangement $\mathcal{A}$. Hence $\mathcal{CL}$ has $n_3$ triple points and at least $10-n_3$ double points. It follows that $n_2+n_3 \geq 10$, a contradiction with \eqref{E11a2}.

\end{proof}

\end{document}
